\section{Number of runs for error reduction}\label{appendix:runs-error-red}

The question is how many runs of $\TF_n$ are needed, given that it answers correctly with probability $c > 1/2$, in order to obtain the correct answer with probability at least $d$ when computing the majority vote.
We estimate the number of runs $r$ using Chernoff's bound.

\begin{theorem}[Chernoff's bound]
    Let $y_1, \dots, y_r$ be independent Boolean random variables with $\Pr[y_i = 1] = c'$.
    Let $Y = \sum_{i=1}^r y_i$ and let $\delta \in (0,1)$.
    Then $\Pr[Y \leq (1-\delta)rc'] \leq e^{-\frac{\delta^2}{4}rc'}$ and $\Pr[Y \geq (1+\delta)rc'] \leq e^{-\frac{\delta^2}{4}rc'}$.
\end{theorem}

Let $y_i$ be a random Boolean variable representing whether the $i$-th run of $\TF_n$ answered correctly on input $x$, i.e., $\TF_n^*(x) = \gL(x)$.
By definition of the family $\{\TF_n\}_{n\in\mathbb{N}}$, $\Pr[y_i = 1] > c > 1/2$.
Therefore, there exists $c \leq c' \leq 1$ such that $\Pr[y_i = 1] = c'$.

We need to find $r$ such that $\Pr[Y < r/2] < d$.
We choose $\delta$ such that $\delta > \frac{2c'-1}{2c'}$; then $r/2 < (1-\delta)rc'$.
Thus, $\Pr[Y < r/2] < \Pr[Y < (1-\delta)rc'] < e^{-\frac{\delta^2}{4}rc'}$.
Finally, choosing $r \leq \frac{4\ln(d)}{c'\delta^2}$, we have $d \leq e^{-\frac{\delta^2}{4}rc'}$, and therefore, $\Pr[Y < r/2] < d$.

Thus, doing $\rceil\frac{4\ln(d)}{c'\delta^2}\lceil$ runs of $\TF_n^*(x)$ and computing the majority vote of the runs is enough to give the correct answer for $\gL(x)$ with probability at least $d$.
Observe that computing more runs only increases the probability of answering correctly.
Thus, for having the right answer with probability at least $d$ for all inputs of size $n$, we run $\TF_n$ for the largest $r$ needed.
Observe that $r$ is largest for the inputs $x$ whose probability of being answered correctly is closest to $c$, i.e., those with the highest probability of error.